\section{瞬时变速首先改变什么}
\subsection{同一位置上的两组守恒量}

火箭短时间点火可近似为位置不变、速度发生有限改变。点火期间有外加推力，机械能和关于中心的角动量一般改变；点火前后的无推力飞行段，才分别保持各自的能量和角动量。应在点火点画出 $\bm v_-$、$\Delta\bm v$、$\bm v_+=\bm v_-+\Delta\bm v$，再计算新的轨道。

单位质量能量的变化为 $\Delta\varepsilon=\bm v_-\cdot\Delta\bm v+|\Delta\bm v|^2/2$。同样大小的顺行脉冲在原速度较大处带来较多能量增加；逆行脉冲则减少能量。若脉冲纯径向，它不改变角动量，却仍改变能量；若脉冲纯切向，二者一般都改变。

点火后先求 $\varepsilon_+=v_+^2/2-\mu/r$、$h_+=|\bm r\times\bm v_+|$，再求 $p=h_+^2/\mu$、$e^2=1+2\varepsilon_+h_+^2/\mu^2$。一个代数上存在的近地点不保证飞船未来会走到那里：还要看当前径向速度和轨道是否闭合。向外的抛物线或双曲线只会越走越远，不能沿时间正方向返回过去的近地点。

\typeout{ASTROFIG:fig:at3:impulse}\fig{fig:at3:impulse}{脉冲与原速度相加后，新的径向分支也要保留}{\begin{tikzpicture}[scale=.85]
\fill(0,0)circle(2pt)node[below]{$F$};\fill(3,0)circle(2pt)node[below]{$A$};\draw[flow](3,0)--(3,1.6)node[left]{$\bm v_-$};\draw[force](3,1.6)--(4.2,1.6)node[above]{$\Delta\bm v$};\draw[flow,blue](3,0)--(4.2,1.6)node[right]{$\bm v_+$};
\end{tikzpicture}}
图\ref{fig:at3:impulse}的径向脉冲使飞船开始向外运动，切向分量保持原值。图中斜边才是点火后总速度，不能把径向脉冲的大小直接当成新轨道速率。\subsection{在近远地点作切向脉冲}

近远地点的径向速度本来为零，切向脉冲后仍为零，该点仍是新轨道的一个端点。但它可能由远地点变成近地点，取决于新速度与当地圆轨道速度的比较。低于圆速时该点是远地点，高于圆速且低于逃逸速度时是近地点，等于圆速时得到圆。

若保持近地点 $r_p$、改变远地点 $r_a$，应在近地点加减速，端点速度由 $v_p^2=2\mu r_a/[r_p(r_p+r_a)]$ 求得。若保持远地点，则在远地点点火。不能把“加速升轨”理解成点火点半径立即增加：瞬间位置不变，改变的是随后到达的另一端。


\section{双切椭圆转移与能量}
\subsection{两次脉冲分别完成不同任务}

从圆轨道 $r_1$ 转到同平面外圆轨道 $r_2>r_1$，选择同时与两圆相切的椭圆，半长轴 $a_t=(r_1+r_2)/2$。起点圆速 $v_{c1}=\sqrt{\mu/r_1}$，转移椭圆近端速度 $v_{t1}=\sqrt{2\mu r_2/[r_1(r_1+r_2)]}$；第一次脉冲为 $v_{t1}-v_{c1}$。

到达外圆时，椭圆远端速度 $v_{t2}=\sqrt{2\mu r_1/[r_2(r_1+r_2)]}$ 小于外圆圆速 $v_{c2}$，还需第二次顺行脉冲 $v_{c2}-v_{t2}$，否则飞船会沿椭圆返回。总速度增量为两次脉冲大小之和，而不是最后速度减去最初速度的标量差。

\typeout{ASTROFIG:fig:at3:transfer}\fig{fig:at3:transfer}{第一次点火进入转移椭圆，第二次点火完成圆化}{\begin{tikzpicture}[scale=.85]
\draw[guide](0,0)circle(1.5);\draw[guide](0,0)circle(3);\draw[world](-0.75,0)ellipse(2.25 and 2.1213203435596424);\fill(0,0)circle(2pt)node[below]{$F$};\fill(1.5,0)circle(2pt)node[right]{$A$};\fill(-3,0)circle(2pt)node[left]{$B$};\draw[flow](1.6,0)--(1.6,1)node[right]{$v_{t1}$};\draw[flow](-3,0)--(-3,-1)node[left]{$v_{t2}$};
\end{tikzpicture}}
图\ref{fig:at3:transfer}的 $A,B$ 相隔半圈，转移时间是椭圆周期的一半。第二次脉冲位于 $B$，并非沿转移途中持续推力。\subsection{轨道能量、推力做功与燃料}

相同卫星质量、同一中心天体的两条圆轨道，总机械能之差为 $\Delta E=\mu m(1/r_1-1/r_2)/2$。在理想无损推力模型中，推力对卫星做功等于这个差；若转移包含大气阻力或非弹性碰撞，推力做功还要补偿耗散。

燃料消耗不能单凭 $\Delta E$ 判定。短时喷气、恒定相对喷气速度 $u$ 的火箭满足 $\dd v=-u\dd m/m$，积分得 $m_f/m_i=\exp(-\Delta v/u)$。多个短脉冲的质量比相乘，对给定 $u$ 等价于把各次 $\Delta v$ 相加。这里的质量变化来自火箭自身携带的喷出物；轨道速度公式用该段飞船质心的速度，不能再把固定质量动能差直接称作全部燃料释放能。

双切椭圆是本讲明确研究的两脉冲方案。对任意半径比、任意多次脉冲、任意飞行时间，不能未经证明宣称它总是全局最省燃料；比较方案时必须统一允许的脉冲数、轨道平面和端点条件。


\section{相遇还需要相位条件}
\subsection{转移轨道相切不等于碰到目标}

地球和外行星若在同平面顺向圆轨道，探测器从地球公转轨道起沿外切转移椭圆走半周。设转移时间 $\tau=\pi\sqrt{a_t^3/\mu_\odot}$，目标角速度 $n_2$，出发时目标相对探测器领先角为 $\phi$。到达时探测器在出发方向的对面，因此 $\phi+n_2\tau=\pi+2k\pi$。

将 $\phi$ 取一个周期内的相位，即 $\phi_*=(\pi-n_2\tau)\bmod2\pi$。若当前领先角为 $\phi_0$、内轨道角速度 $n_1>n_2$，等待时相位按 $\phi(t)=\phi_0-(n_1-n_2)t$ 改变。求第一次非负等待时间时，还需加适当整数周，不能得到负时间就直接取绝对值。

相位计算默认出发前探测器已在内圆轨道稳定运行，而且火箭点火时间足够短。若题目给的是地面发射时刻，脱离地球的有限飞行时间和发射方位也要计入；模型中被忽略的过程不能悄悄混入一个日历答案。

\subsection{共转位置的引力与转动效应}

日地连线上共转点的轨道角速度必须与地球相同，因而它不是“日地引力恰好抵消”的位置。在随日地线转动的坐标系可写有效力为零，在惯性系则写引力合力提供对应向心加速度，两者等价。

靠近小天体的共线点，距离量级为 $R(m/(3M))^{1/3}$。这个三次方关系来自大天体引力梯度、共转向心加速度变化与小天体平方反比吸引之间的平衡。系数三有明确的参考系和几何含义，不是所有“潮汐力”问题都可照搬的系数。


\section{到达边界与最小发射速度}
\subsection{先判断允许的轨道弧段}

从球面一个点到另一个点，忽略大气和自转时，可以使用一段与球面相交的椭圆弧。完整数学椭圆的部分轨迹会穿过球体，但飞行只取位于表面以上、起终点之间的那段弧。讨论最小速度时必须同时满足两个端点、轨迹在外侧和初始朝外等条件。

对称的端点等半径问题，可利用椭圆第二焦点或参数极值求最优半长轴。固定起点半径时，速度平方 $v^2=\mu(2/R-1/a)$ 随半长轴增加而增加，所以最小速度对应最小可行的 $a$。这是一条受几何约束的最值问题，不是简单地让总能量尽可能负；若不保留端点几何，最小能量会给根本到不了目标的轨道。

无推力飞行中，角动量和能量检查的是整条轨道；可达性还要检查实际选取的时间方向。撞击条件通常是最近距离不大于表面半径，而安全掠过取严格大于。恰好相切时的等号能否允许，应在题设中明确。


\section{变轨方案的比较条件}
\subsection{速度方向也消耗脉冲}

  若变轨前后速率为 $v_1,v_2$，方向夹角为 $i$，所需脉冲是两矢量之差，$\Delta v=\sqrt{v_1^2+v_2^2-2v_1v_2\cos i}$。只改变轨道平面而不改速率时，$\Delta v=2v\sin(i/2)$，机械能不变但仍消耗燃料。把能量差为零误认为不用点火，会漏掉方向改变。

  在给定交点允许改变平面的条件下，低速处通常更省方向脉冲。但如果为此额外抬高远地点，抬轨和恢复轨道的脉冲也要计入；单独比较某一点的平面改变成本，不是完整方案的总成本。与此同时改变速率和方向时，两种脉冲应作矢量合成，不能把两个标量成本简单相加后宣称最优。
  
\subsection{慢推力与瞬时推力的差别}

  在近圆的缓慢螺旋升轨中，推力持续沿速度方向，轨道随时间变化；圆轨道关系可作为缓变近似。单位质量能量 $\varepsilon=-\mu/(2r)$，推力功率为 $a_tv$，因此累计推力速度增量满足 $\dd V_{\rm thrust}=\dd\varepsilon/v=\sqrt\mu\,\dd r/(2r^{3/2})$。积分得到 $\Delta V_{\rm thrust}=\sqrt{\mu/r_0}-\sqrt{\mu/r_1}$。

  当目标半径趋无穷远，这个理想慢推力极限为初始圆速，而一次瞬时逃逸只需其 $\sqrt2-1$ 倍。它们端点能量差可以相同，但推力在不同位置、不同速度下做功，累计速度增量不同。比较时应同时列出端点、允许的飞行时间和推力形式。
  
\subsection{相位轨道与中间轨道的安全边界}

  为在同一圆轨道上追赶目标，可暂时进入周期更短的内侧椭圆，绕行后在原点再次圆化。相位改变来自周期差，不能只让瞬时速度更大就认定平均角速度更大：在原圆上切向减速反而进入更短周期的内轨道。

  任何内侧相位轨道都要检查近日点是否高于实体表面或大气；外侧相位轨道要检查远地点与影响区边界。对于多次脉冲方案，还要逐段计算飞行时间。把中间半径取无穷远得到的速度增量极限，通常对应无穷长转移时间，不能说有限时间已经达到这个极限。
  

\section{喷气质量变化与最值的检查}
\subsection{火箭方程的符号和系统边界}

  取短时间内火箭与即将喷出的燃料为物质系统。火箭原质量 $m$、速度 $v$；喷出后火箭质量为 $m+\dd m$，其中 $\dd m<0$，速度为 $v+\dd v$。气体质量是 $-\dd m>0$，相对火箭向后速度为 $u$，其惯性系速度在一阶近似下为 $v-u$。

  忽略外力短时冲量，初末动量相等：
  \[mv=(m+\dd m)(v+\dd v)+(-\dd m)(v-u).
  \]
  展开后 $v\dd m$ 项抵消，舍去二阶小量 $\dd m\dd v$，得到 $m\dd v=-u\dd m$。若喷速恒定且喷气方向始终与所需脉冲一致，积分为 $\Delta v=u\ln(m_i/m_f)$。此处的负号来自质量减少，与“发动机方向向后”是同一个物理事实的两种记法，不能重复添加负号。
  
\typeout{ASTROFIG:fig:at3:rocket}\fig{fig:at3:rocket}{取火箭和即将喷出部分为同一短时系统}{\begin{tikzpicture}[scale=.85]
\draw[world](0,0)rectangle(2,1);\node at(1,.5){$m,v$};\draw[flow](.4,1.5)--(1.8,1.5);\draw[world](4,0)rectangle(5.7,1);\node at(4.85,.5){$m+\dd m$};\draw[world](2.8,.2)rectangle(3.5,.8);\node[above]at(3.15,.8){$-\dd m$};\draw[force](3.3,-.2)--(2.3,-.2)node[below]{相对喷速 $u$};
\end{tikzpicture}}

  图\ref{fig:at3:rocket}中，喷出物仍属于初末动量比较的物质系统。若只对剩下的火箭写 $F=\dd(mv)/\dd t$ 而不说明流出动量，就会遗漏或重复质量流项。轨道题常把这个过程近似成瞬时速度跳变，但推导质量比时仍需尊重变质量系统的边界。

  多次脉冲时，第一个脉冲的末质量成为第二个脉冲的初质量，所以总质量比为各段指数因子的乘积，最终指数中出现各次速度增量大小之和。若不同发动机的喷速不同，则应分别累计 $\Delta v_j/u_j$，而不能再只按总速度增量比较燃料。若点火间有抛弃整级或质量补给，初末质量也须按阶段更新。

  求某方案“最省”的时候，要先固定比较集合：同一初末轨道是否还要求同一轨道平面、同一到达相位、相同飞行时间上限，以及中间轨道是否可以无限抬高。改变其中一项，最值问题就改变了。双切、双椭圆、慢推力的比较正是用不同可行集合说明这一点，不能混成一个没有约束的口号。

  从解答规范看，每个脉冲应给出前后速率、方向与所在点，每段飞行应给出新的能量、角动量或端点条件。若以某段轨道的最高点作为下一次点火点，还应确认那是未来真正会到达的位置。算出一个负的等待时间、负的近日点或高于允许表面的掠近条件，都应回到模型分支检查，不能通过取绝对值或删去符号把它变成答案。
  
